% 2 Background and related work -- 0.9 page
\section{Background and Related Work}

\subsection{Kernel Tracing with LTTng and Trace Compass}

LTTng~\cite{desnoyers2006lttng} records Linux kernel events with low overhead: system call entries and exits, scheduler switches and wake-ups, timers, interrupts, and block-device, network and signal events; Gebai and Dagenais~\cite{gebai2018survey} survey the Linux kernel and user-space tracers and their overhead, and kernel-level events have been shown to enrich empirical performance research~\cite{noferesti2024kernel}. Trace Compass~\cite{tracecompass} turns these events into \emph{state systems}, time-indexed models of each thread's state, current system call and each CPU's run queue, and into an \emph{execution graph} of which thread woke which. The \emph{critical path} of a thread follows this graph backwards through its waits to the work that delayed it~\cite{giraldeau2016wait}. Critical paths also profile latency at scale in distributed services~\cite{eaton2022cpt,zhang2022crisp}, and Trace Compass's critical path has been used to detect anomalies with an LSTM at a fraction of the cost of system-call features~\cite{nourollahi2025tracelens} and to find performance archetypes and regressions~\cite{shahedi2026archetypes}. These analyses are shown in views such as Control Flow, Resources and Critical Path. The headless Trace Compass Trace Server exposes the same analyses through the Trace Server Protocol~\cite{tsp}, which TMLL uses. Reading these views well requires knowing what each thread state means and where to look, which is the expertise \sysname aims to supply.

\subsection{Related Work}

\textbf{Language models over traces.} TAAF~\cite{ezaz2026taaf} answers natural-language questions about LTTng kernel traces by building a time-indexed Trace Compass state system and, for each question, extracting a knowledge graph of threads, CPUs and resources over which an LLM reasons. On TraceQA-100, 100 questions about slices of a SciMark~2.0 trace, it improves accuracy by up to 31.2\% over the state system's numeric output, mostly on multi-hop and causal questions; our hub-scope labels follow its distinction between single- and multi-hop questions. \sysname differs in three ways. Its analyzers compute the measurements an answer needs, such as wait chains, per-round finish times and stuck waits, so the LLM explains results rather than deriving them from a graph. It analyzes whole traces rather than slices. And our evaluation asks each question of traces with and without the symptom, so it measures false positives.

\textbf{Machine learning for Trace Compass.} TMLL~\cite{shahedi2025tmll}, built with Ericsson, applies anomaly detection, change-point detection, correlation and capacity planning to the Trace Server's analyses. Its authors report that technical excellence alone did not drive adoption by practitioners; explainable results and expert knowledge built into the interface did. \sysname uses TMLL as one source of supporting evidence, and we evaluate TMLL's analyses, given to an LLM, as a baseline.

\textbf{Deterministic evidence with LLM explanation.} Outside kernel tracing, two systems share \sysname's division of labour. RCACopilot~\cite{chen2024rcacopilot} matches each cloud incident at Microsoft to a handler that collects diagnostic data, then has an LLM predict the root-cause category and explain it, reaching an accuracy of up to 0.766 on a year of incidents. Opal~\cite{zaeed2025opal} adds hardware-counter and Roofline analytics to LLM prompts to optimize GPU kernels; even a single source of insight produced a speedup in over 98.5\% of 1{,}640 experiments, and its one incorrect result stemmed from a vague diagnostic, which supports giving the LLM precise findings. LLMs have also been fine-tuned to recommend root causes and mitigations for over 40{,}000 cloud incidents~\cite{ahmed2023recommending}, and Praxium diagnoses cloud anomalies from telemetry and service dependencies~\cite{kumar2026praxium}. Closer to kernel traces, Fournier et al.~\cite{fournier2021cause} find the causes of slow web requests by clustering and comparing their kernel-level executions. \sysname brings the division of labour to kernel traces, where the evidence must first be computed from millions of events, and to open questions rather than classification or code generation.

\textbf{Coding agents.} General coding agents such as Codex~\cite{codexcli} interleave reasoning with tool use~\cite{yao2023react} and are evaluated mostly on resolving code issues~\cite{jimenez2024swebench}; given a shell, they can run tools over a trace and compute their own evidence. We include one as a baseline: it shows what an engineer gets today without a trace-specific assistant, and at what cost.
